\subsection{The actual factorial free source and square frame}
The selected source is the literal 239-coordinate Strict Horizons source,
not the factorial-authenticating Pooled Powers source. Its semantic values
are
\[
 t=16+h^2,\quad F=t^t,\quad B=\binom{F+t}{t},\quad
 q=4z+1,\quad H=2+C+\sum_{j=1}^8A_j.
\]
Here $B$ is initially only a proposed binomial value. There are nine actual
CRT representatives, $C,A_1,\ldots,A_8$, and a globally evaluated derived
label $J(\boldsymbol A,C)$, which is not an additional independent column.
The unchanged endpoints and index row are
\begin{equation}\label{eq:strict-endpoints}
 x=F+1+(y+1)u_0,\qquad
 x=1+[1+(t+1)y]u_t,\qquad Bh_B=C+F+1.
\end{equation}
The separate request at the already positive base $t$ will first recover
$F=t^t$. No coordinate is authenticated as $t!$ in this branch.

Before using the signs of these arithmetic values, the retained squared
semantic anchor and deterministic comparison-center proof give all
fourteen canonical signed Pell pairs. The literal actual-atom frame is
$W=2+\sum a^2$, with exactly 34 squares. Its atoms comprise
$t,F+t,x,y+1,F,B$; the nine representatives; $N$ and the global derived
label; $q,q+1,BH$; and the fourteen remaining request outputs.
Consequently every base, index and output of all seventeen requests has
square below $W$, and $t^2$ gives $W\ge258$.
The frame provides $|A|>W^W$ independently of the later sign bootstrap.
Stored coordinate offsets matter: for example the source coordinate
named $y$ stores $y+1$, and one schema coordinate stores $A_a+1$.
The checker expands the actual affine atoms, not their historical names.

Apply Theorem~\ref{thm:pooled-exactness} to eleven requests at $q+1$
and three at $q$, leaving the requests at $t,B,BH$ separate.
Each group replaces $r$ private quotients and one modulus by two
coordinates, saving $r-1$ variables and rows: ten plus two.
The retained $AY_i$ wires make every new row quadratic.

\subsection{An odd beta multiplier and a wider sphere}
Replace the old $y=F(2v+1)^2$ by
\begin{equation}\label{eq:odd-beta}
 y=F(2v+1).
\end{equation}
The multiplier's defining row becomes affine and permits its monic
elimination. This is a semantic extension followed by an exact projection,
not a claim that every new multiplier is an old square.

For the eleven normalized remainders retain the constant one and replace
coefficient four by two:
\begin{equation}\label{eq:strict-sphere}
 M^2=1+2\sum_{j=1}^{11}\rho_j^2+G(a,b,c),
 \qquad G(a,b,c)=a^2+b^2+c(c+1).
\end{equation}
For integers $G\ge0$, so $M\ne0$ and
$|\rho_j|<|M|/\sqrt2<3|M|/4$.
The constant one is essential to this nonzeroness argument and is not
removed. The old signed-factorial proof of nonzeroness is not imported
into this factorial-free source.

For a normalized row $Dv=M$, $vU=MQ+\rho$, the nonzero denominator
and $M$ give $v\ne0$. The row gives $v\mid\rho$, so $r=\rho/v$
is integral and
\begin{equation}\label{eq:strict-remainder}
 U=DQ+r,\qquad |r|<3|D|/4.
\end{equation}
No sign of $D,v,M$ or $\rho$ was assumed.

\subsection{Recover the signs in an acyclic order}
The unpooled $t$ request first gives $F=t^t>2^{32}$.
Equation~\eqref{eq:odd-beta} implies $|y|\ge F>1$, but does
\emph{not} yet imply $y>0$. The affine radix has $q\ne0$ and
$|q+1|\ge2$. Apply the $q+1$ pool: all nine actual representatives
and the derived-label index are nonnegative and all their numerators
are exact. Hence $H\ge2$ and the index row makes $B\ne0$.
Now $|BH|\ge2$, so its separate request gives $x\ge0$ and
$q=(BH)^x$.

If $x=0$, the last endpoint in \eqref{eq:strict-endpoints} would
require $1+(t+1)y$ to divide $-1$. Its absolute value is at least
$(t+1)F-1>1$, impossible. Thus $x\ge1$ and $|q|\ge2$.
Only now use the $q$ pool to recover exact $q^t,q^{y+1},q^x$
and nonnegative indices. In particular $y\ge-1$; together with
$|y|\ge F>1$, this gives $y\ge F>0$. This is the first justified
use of the positive sign of $y$. The ratio $q^tv=q^x$ also gives
$x\ge t$.

If $u_0\le-1$, the first endpoint would give $x\le F-y\le0$.
Therefore $u_0\ge0$, $x\ge F+1>1$, $u_t\ge1$ and
\begin{equation}\label{eq:strict-growth}
 x\ge2+(t+1)y\ge2+(t+1)F.
\end{equation}
Writing $Q=|q|$ and $b=|B|$, we have
$Q=|BH|^x\ge2^x$ and $b<Q/2$.
No positivity of $B$ or $q$ was used to obtain these bounds.

\subsection{Decode the binomial at a signed radix}
Let $C_j=\binom{F+t}{j}$. For $0\le j\le t$,
\[
 C_j\le(F+t)^t\le(2F)^t\le2^{t(F-1)}.
\]
Here $t\le F$ and $2F\le2^{F-1}$ for $F\ge4$.
Equation~\eqref{eq:strict-growth} implies
$C_t<Q/32$ and $\max_{j\le t}C_j/(Q-1)<1/32$.
The signed low tail $r_0=\sum_{j<t}C_jq^j$ therefore satisfies
$|r_0|<Q^t/32$, even for negative $q$.
The restored carrier-digit row is
\[
 (q+1)^{F+t}=q^t(B+q h_{\rm high})+r.
\]
Thus $q^t\mid r-r_0$, while
$|r-r_0|<(3/4+1/32)Q^t<Q^t$ by
\eqref{eq:strict-remainder}. Hence $r=r_0$.
Subtract and divide by $q^t$ to obtain $B\equiv C_t\pmod q$.
Since $|B-C_t|<Q/2+Q/32<Q$, we conclude
\begin{equation}\label{eq:strict-binomial}
 B=\binom{F+t}{t}>0,\qquad q=(BH)^x>0.
\end{equation}
These signs are consequences of the decoded digit, not its premises.

For every other genuine low digit, the retained source estimate is
$0\le r_0<D/32$. The same divisible-difference proof with radius
$3D/4$ gives $r=r_0$. Thus the old strict interval witnesses
$0\le r<D$ can all be restored, not merely the necessary congruences.
Ordinary upper representatives have binomial coefficients at the required
positions at most $H^x$. The global derived label has bound
$J(\boldsymbol A,C)\le H^{2^{22}}$, with $2^{22}t<x$.
Since $q=B^xH^x$, these are the retained margins giving the $1/32$
tails. The source's local majorant uses exponent $2^{23}$;
$2+t+x+y+F<2x$ and $x>2^{32}$ give
\[
 (2+t+x+y+F)^{2^{23}}<(2x)^{2^{23}}<2^x.
\]
The last inequality follows from monotonicity of
$x/\log_2(2x)$ above $2^{32}$ and its value there.

\subsection{The retained prime power and flow interfaces}
We record why the now restored source conditions suffice.
Let $\Lambda_t=\operatorname{lcm}(1,\ldots,t)$.
The elementary bound $\Lambda_t\le4^{t-1}$ follows by induction:
for $k=\lceil t/2\rceil$, each additional prime-power factor in
$\Lambda_t/\Lambda_k$ divides the central binomial coefficient,
which is at most $2^{t-1}$. Therefore
$\Lambda_t^2\le16^{t-1}<t^t=F$ for $t\ge16$.
This licenses the linear CRT radius $F$, without a factorial coordinate.

Two different prime-power facts are needed. First, if
$p^e\Vert\binom{F+t}{t}$, then $p^e\mid F+i$ for some
$1\le i\le t$. Indeed the factorial valuation is a sum of zero-one
terms at prime-power levels, and is at most
$\max_i v_p(F+i)$; beyond that maximum no level contributes.
Second, every row $i$ has such a full prime power $p^e>t$.
To see this, $F+i>\Lambda_t^2$ gives a prime with
$v_p(F+i)>2v_p(\Lambda_t)$. All other interval terms have
$v_p(F+j)=v_p(j-i)$. The loss on division by $t!$ is
\[
 v_p\left(t\binom{t-1}{i-1}\right)\le v_p(\Lambda_t),
\]
since each Legendre summand is zero or one and vanishes above $t$.
The remaining full exponent is larger than $v_p(\Lambda_t)$, so
$p^e>t$ and still divides $F+i$.

For $A\in\N$, divisibility $p^e\mid\binom Ak$ forces
$A\bmod p^e<k$. If $A\ge k$ and its falling interval contained
no multiple of $p^e$, all valuation levels from $e$ onward would
vanish; the lower zero-one levels sum to at most $e-1$, a contradiction.
If $A<k$, the conclusion is immediate. A row's large selector decodes
its bounded local tuple; the restored strict ranges identify the row,
and the local bound $0\le\mathcal L<2^x\le p^{xe}$ makes its
divisible local error zero. The derived label is decoded by polynomial
congruence from the same tuple, not treated as an independent column.
For the converse CRT construction, the every-prime fact assigns a
compatible valid row at \emph{every} prime of $B$. Precision $p^{2xe}$
loses fewer than $x\le xe$ factors on division by each relevant factorial,
leaving the required $p^{xe}$ divisibility. One of these two prime-power
lemmas alone would not establish both witness directions.

If $y=F$, the initial beta accumulator is zero. Positive-weight local
updates propagate zero, contradicting the final accumulator one.
Thus $y>F$, and the restored updates give
\[
 \prod_{i<t}(1+Fv_i)=(1+F)\prod_{i<t}(1+Fp_i)(1+Fq_i).
\]
Discard zero labels. Unequal factor counts would have a smaller count
$r\le t$ and force $(F+1)^{r+1}\le t^r(F+1)^r$, impossible
since $F+1>t^t$. Hence both products have the same number $r\le t$
of positive factors. In their product polynomials every coefficient of
positive degree lies in $[1,F]$, so coefficient differences have absolute
value below $F$. Reduction modulo $F$ at the first nonzero difference
coefficient shows that equality at $F$ implies polynomial equality.
Distinct factors $1+jz$ then recover exactly the required label
multiplicities. An individual coefficient is allowed to equal $F$;
the strict bound is on its difference.

Finally, take the finite Horn closure of truth. If falsity were outside,
every listed rule with an outside conclusion would have an outside
premise. Summing conclusions minus premises over outside labels would
be at most zero; the recovered multiplicities make it one. This
contradiction gives a C222 refutation. The identified C222/ZF/ZFC
interfaces in Section~\ref{sec:foundation} complete soundness.

\subsection{Constructive completeness and retained choice order}
Begin with a finite calibrated C222 refutation. The unchanged Strict
Horizons construction chooses $t=16+h^2$ strictly beyond its finite
tree length, local parameters and derived labels, with $h$ even.
It then forms $F=t^t$, the finite pending-accumulator path and an old
spacing $y=F(2v_{\rm old}+1)^2$ large enough for that path.
The even choice gives $v_2(F)=t\,v_2(t)>v_2(t!)$;
a sufficiently large odd multiple of the odd part of $t!$ supplies
the old multiplier. The beta moduli are pairwise coprime and odd,
so their CRT code $x$ can be chosen even and large. The retained
binomial CRT at $N^2=B^{2x}$ gives all representatives.
The binomial $B$ is odd because $v_2(F+i)=v_2(i)$ for $i\le t$.
Thus parity-adjustable representatives supply an even dominant derived
anchor and odd $H$. Consequently $(BH)^x\equiv1\pmod4$,
which supplies the new affine radix coordinate.
Only after the arithmetic atoms are fixed are fresh simultaneous
canonical Pell/frame witnesses selected.

Keep all these old source values and set
\[
 v_{\rm new}=\bigl((2v_{\rm old}+1)^2-1\bigr)/2\in\Z.
\]
This preserves actual $y$ and every frame atom under
\eqref{eq:odd-beta}. The true normalized remainders satisfy
$|\rho_j|<M/32$, so
$2\sum_j\rho_j^2<22M^2/1024<M^2$.
The positive integer $M^2-2\sum\rho_j^2$ is at least one,
therefore $M^2-1-2\sum\rho_j^2\ge0$.
Every nonnegative integer $g$ is represented by $G$: the three-square
theorem applied to $4g+1$ gives exactly one odd square modulo four,
which yields $g=a^2+b^2+c(c+1)$ on division by four.
This supplies the wider sphere. Lemma~\ref{lem:pooled-identity}
supplies the four pooled coordinates. All new rows now vanish.
This is a constructive existence equivalence, not a bijection between
arbitrary old and new beta multipliers or ternary-guard triples.

\subsection{The exact algebra after the semantic changes}
The private-wire reduction first stores $J_0=s_0^2+g^2$ instead of
two coordinates used only through their sum. Restoration gives the two
original squares. It next stores $J=FS$, retaining $K=FS'$, instead
of three private factors $1+Fb$, $1+F(a-1)$ and $1+Fd$,
where $d=6a+8b+2p$. The consumer identities are
\[
 S'(1+Fb)=S'+Kb,\quad S(1+F(a-1))=S+J(a-1),
\]
\[
 S(1+Fb)-S'(1+Fd)=S+Jb-S'-Kd.
\]
Every old factor restores by its defining polynomial; the certificate
checks the residual differences as multiples of the retained monic
product definitions. The sum and product replacements save three
variables and rows without creating cubic residuals.

Next put $D_i=X_i-Z_i$, with inverse $X_i=D_i+Z_i$.
Subtract the fourteen retained norm rows from the orientation row.
The pooled second row now uses $(m_i-D_i)^2-Y_i^2$.
These invertible shears and row operations change density, not dimension.
Every transported Round~20 core operation is rechecked on this actual
source: either an integral shear, a reversible row operation, or a monic
projection with an integral degree-at-most-two restoration. A quadratic
restoration is allowed only when every other occurrence of the deleted
coordinate is linear. The new low-height core exposes two further
common-Hessian directions, handled by
Theorem~\ref{thm:radical-projection}. The dimension companion uses one
additional affine consequence as in Lemma~\ref{lem:affine-elimination}.
Finally frozen integer row-basis operations select the support/height
tradeoffs. No operation is licensed merely because its old coordinate
name occurs in an earlier certificate.
